% ============================================================
% CHƯƠNG 3: THIẾT KẾ VÀ XÂY DỰNG HỆ THỐNG
% ============================================================
\chapter{Thiết kế và xây dựng hệ thống}

\section{Phân tích yêu cầu}

\subsection{Yêu cầu chức năng}
% TODO:
% FR1: Tóm tắt tự động bài viết tin tức tiếng Việt
% FR2: Kiểm chứng sự kiện dựa trên highlight text
% FR3: Hỗ trợ đa mô hình ngôn ngữ (multi-model)
% FR4: Tự động định tuyến chọn model phù hợp (routing)
% FR5: Chế độ đánh giá so sánh đa model (evaluation mode)
% FR6: Dashboard theo dõi hoạt động người dùng
% FR7: Trang metrics hiển thị kết quả đánh giá

\subsection{Yêu cầu phi chức năng}
% TODO:
% NFR1: Thời gian phản hồi < 10 giây (streaming)
% NFR2: Hoạt động trên Chrome, Edge, Brave
% NFR3: Không ảnh hưởng hiệu năng trang web gốc (Shadow DOM)
% NFR4: Bảo mật API key (server-side only)

\section{Kiến trúc hệ thống}

\subsection{Tổng quan kiến trúc}
% TODO:
% - Kiến trúc microservice 3 thành phần
% - Sơ đồ kiến trúc tổng quan
% Hình 3.1: Sơ đồ kiến trúc hệ thống Fiber

\subsection{Tiện ích mở rộng trình duyệt (Browser Extension)}
% TODO:
% - Plasmo framework, React 18, TypeScript, Tailwind CSS
% - Content Script injection vào các trang báo
% - Shadow DOM isolation
% - Sidebar (tóm tắt) và Popover (kiểm chứng)
% Hình 3.2: Kiến trúc browser extension

\subsection{Máy chủ API (Backend Server)}
% TODO:
% - Next.js 14 App Router
% - API routes: /api/summarize, /api/fact-check, /api/metrics, /api/evaluate
% - Services layer: llm, summarize, fact-check, evaluation, routing
% - Zod schema validation cho request/response
% Hình 3.3: Sơ đồ các API endpoint

\subsection{Dịch vụ BERTScore (BERTScore Microservice)}
% TODO:
% - FastAPI, Python, PhoBERT (vinai/phobert-base)
% - POST /calculate-score
% - Truncation 256 tokens
% - Triển khai trên HuggingFace Spaces
% Hình 3.4: Kiến trúc BERTScore microservice

\section{Thiết kế chi tiết các module}

\subsection{Module tóm tắt văn bản}
% TODO:
% - Luồng xử lý: URL → Content Extraction → Prompt → LLM → Structured Output
% - Prompt engineering cho tóm tắt tiếng Việt
% - Structured output với Zod schema
% - Streaming response (SSE)
% Hình 3.5: Sequence diagram - Quy trình tóm tắt

\subsection{Module kiểm chứng sự kiện}
% TODO:
% - Luồng: Selected text → Tavily Search → Context building → LLM reasoning → Trust score
% - Thiết kế prompt cho fact-checking
% - Trust scoring system (High/Medium/Low)
% Hình 3.6: Sequence diagram - Quy trình kiểm chứng

\subsection{Module đa mô hình (Multi-Provider LLM)}
% TODO:
% - Kiến trúc provider abstraction (OpenAI, Google Gemini, Anthropic Claude)
% - 14 models across 3 providers
% - Xử lý đặc thù từng provider (temperature, structured output, streaming)
% - Model configuration trong database
% Bảng 3.1: Danh sách các mô hình được hỗ trợ

\subsection{Cơ chế định tuyến mô hình (Routing Mechanism)}
% TODO:
% - Bài toán: chọn model phù hợp dựa trên độ phức tạp bài viết
% - Phân loại độ phức tạp: simple (< 500 tokens), medium (500-1500), complex (> 1500)
% - Routing logic: simple → ViT5, medium → PhoGPT/GPT-4o, complex → GPT-4o
% - Output Fusion: kết hợp kết quả từ nhiều model trong evaluation mode
% Hình 3.7: Flowchart cơ chế routing
% Bảng 3.2: Ma trận routing (complexity × model)

\subsection{Module đánh giá (Evaluation Module)}
% TODO:
% - Evaluation mode: chạy đồng thời nhiều model
% - Tính toán ROUGE, BLEU, BERTScore, Compression Rate, Latency
% - BERTScore selection: chọn kết quả tốt nhất
% - Lưu trữ metrics vào Supabase
% Hình 3.8: Sequence diagram - Evaluation mode

\section{Thiết kế cơ sở dữ liệu}

\subsection{Sơ đồ quan hệ}
% TODO:
% - Bảng evaluation_metrics: lưu kết quả đánh giá
% - Bảng dashboard_actions: log hành động người dùng
% - Bảng model_configurations: cấu hình mô hình
% Hình 3.9: ER Diagram

\subsection{Mô tả chi tiết các bảng}
% TODO:
% Bảng 3.3: Cấu trúc bảng evaluation_metrics
% Bảng 3.4: Cấu trúc bảng dashboard_actions

\section{Thiết kế giao diện người dùng}

\subsection{Sidebar tóm tắt}
% TODO: Hình 3.10: Mockup/Screenshot sidebar

\subsection{Popover kiểm chứng}
% TODO: Hình 3.11: Mockup/Screenshot fact-check

\subsection{Trang Settings và Debug}
% TODO: Hình 3.12: Trang settings với model selector

\subsection{Trang Metrics}
% TODO: Hình 3.13: Trang metrics với biểu đồ

\section{Tổng kết chương}
% TODO: Tóm tắt các quyết định thiết kế chính
